\section{Related Work}
\label{sec:related}

\subsection{Uncertainty Quantification via Sampling}

Semantic entropy~\cite{kuhn2023semantic} estimates an LLM's uncertainty
about a free-form generation by sampling multiple outputs, clustering them
by semantic equivalence, and computing the entropy of the resulting cluster
distribution. This requires $K$ generation calls per uncertainty estimate
and is the conceptual basis for a family of clarify-then-generate systems
for code, including ClarifyGPT~\cite{clarifygpt}, which uses repeated
sampling of the coding model's outputs to decide when an ambiguous
requirement warrants a clarifying question. These approaches are effective
but tie the cost of every clarification decision to $K$ additional calls to
the coding model itself, which is frequently the most expensive component
of the pipeline.

\subsection{Single-Forward-Pass Probing}

Kossen et al.~\cite{oatml_sep} introduce semantic entropy probes: a linear
probe trained on a single forward pass' hidden states to predict the
semantic entropy that would otherwise require $K$ generations to estimate,
amortizing the sampling cost of training many probes' worth of labels into
a one-time cost, with cheap inference thereafter. Our probing methodology
directly reuses this approach -- per-layer linear probing of residual-stream
activations, best-split binarization for layer selection, bootstrapped
evaluation -- but applies it to a different prediction target. Where
Kossen et al.\ probe for a property of the model's own \emph{output}
distribution (how much its generations would disagree with each other), we
probe for a property of the \emph{input} prompt (how much unresolved
specification ambiguity it contains) and, in a second stage, a property of
a candidate answer relative to that input (how well it resolves the
ambiguity, for ranking purposes). This is a substantively different
prediction problem: no sampling, clustering, or generation from the model
being probed is involved anywhere in our pipeline, at either training or
inference time, and consequently the model being probed need not be the
model doing the coding.

\subsection{Ambiguity and Clarification in Software Engineering}

Resolving requirements ambiguity before implementation is a long-standing
concern in software engineering, predating LLM-based tooling. Recent
LLM-era work targets this directly for code generation, constructing
benchmarks with deliberately injected or naturally occurring ambiguity and
evaluating whether models recognize the need to ask a clarifying question
before writing code~\cite{clarifycodebench}. Our evaluation builds on
ClarifyCodeBench and on Orchid-style ambiguity injections over
HumanEval~\cite{humaneval} and BigCodeBench~\cite{bigcodebench}, which
construct matched clear/ambiguous prompt pairs by deleting or rewriting a
targeted specification detail and pairing the injected prompt with tests
that specifically exercise the resulting ambiguity.

\subsection{Agentic Coding Benchmarks}

Our live evaluation loop -- eliciting candidate clarifying questions,
ranking them, generating code conditioned on the top-ranked answer, and
scoring the result against real executable test suites -- follows the
general pattern of agentic coding benchmarks that emphasize execution-based
correctness over surface-level similarity to a reference
solution~\cite{humaneval,bigcodebench}. We report correctness using the
fractional test-case pass rate rather than an all-or-nothing pass metric,
and show in Section~\ref{sec:results} that this choice materially affects
which relationships in the data reach statistical significance -- a point
we believe is broadly relevant to how agentic coding pipelines should be
evaluated, independent of our specific method.
